\section{Conclusion}
\label{sec:conclusion}

We interpret on-policy distillation as an RL-like process in which the teacher acts as an implicit reward model for the student's own responses.
Across our experiments on competition mathematics, OPD improves sampling efficiency without expanding the observed solvable set.
Its gains primarily come from making correct responses already accessible to the initial student easier to sample.
This view also helps explain why OPD can collapse despite normal convergence.
In the collapsed setting, the student increasingly follows the teacher's preferences while producing longer and more repetitive responses.
These responses receive high teacher advantage even before training, although the teacher rarely generates them itself.
Masking the loss on truncated responses and initializing the student from an SFT checkpoint mitigate this failure while keeping the teacher fixed.
